\subsection{替换}

设 E 为一个 A-代数。E 上的一个拓扑称为线性拓扑，如果它在平移之下不变，并且存在一个由 E 的理想组成的 0 的基本邻域系（《一般拓扑学》，III，第223页）。此时 E 上的拓扑与其 A-代数结构相容（当 A 赋有离散拓扑时）。带线性拓扑的 A-代数称为线性拓扑化 A-代数。

命题 4. --- 设 I 是一个集合，E 是一个结合、交换、含单位元、线性拓扑化、分离且完备的 A-代数。

\begin{enumerate}
\def\labelenumi{(\roman{enumi})}
\tightlist
\item
  设 \(\varphi\) 是从 A{[}{[}I{]}{]} 到 E 中的连续同态，且 \(x_{i} = \varphi (\mathbf{X}_{i})\) 。那么：
\end{enumerate}

\begin{enumerate}
\def\labelenumi{(\alph{enumi})}
\item
  对所有 \(i \in \mathbf{I}\) ，\(x_{i}^{n}\) 当 \(n\) 趋于 \(+\infty\) 时趋于 0 ；
\item
  若 I 是无限集，则 \(x_{i}\) 沿 I 的有限子集的补集所构成的滤子趋于 0 。
\end{enumerate}

\begin{enumerate}
\def\labelenumi{(\roman{enumi})}
\setcounter{enumi}{1}
\tightlist
\item
  设 \(\mathbf{x} = (x_i)_{i\in I}\) 是 E 中满足 (i) 的条件 \(a\) 与 \(b\) 的元素族。则存在一个且仅一个保单位的连续同态 \(\varphi\) ，它从 A{[}{[}I{]}{]} 到 E 中，使得 \(\varphi (\mathbf{X}_i) = x_i\) 对所有 \(i\in I\) 。
\end{enumerate}

对所有 \(i \in I\) ，\(X_{i}^{n}\) 在 A{[}{[}I{]}{]} 中显然趋于 0 ，当 n 趋于 \(+\infty\) ；另一方面，当 I 为无限集时，\(X_{i}\) 沿 I 的有限子集的补集所构成的滤子趋于 0 。这就证明了 (i)。

设 \((x_{i})_{i\in I}\) 是 E 中满足 (i) 的条件 \(a)\) 与 \(b)\) 的元素族，设 \(\psi\) 为由 \(u\mapsto u((x_i)_{i\in I})\) 给出的、从 \(\mathbf{A}[(\mathbf{X}_i)_{i\in I}]\) 到 E 中的同态，并设 V 是 E 的 0 的一个邻域，且它是 E 的一个理想。由 \(b)\) ，存在 I 的有限子集 J 使得 \(x_{i}\in \mathbf{V}\) 对所有 \(i\in \mathrm{I - J}\) 。其次，由 \(a)\) 存在整数 \(n\geqslant 0\) 使得 \(x_{i}^{n}\in \mathbf{V}\) 对所有 \(i\in \mathrm{J}\) 。设 \(\beta\) 是 \(\mathbf{N}^{(I)}\) 中这样的元素：\(\beta_{i} = n - 1\) 对 \(i\in \mathrm{J}\) ，且 \(\beta_{i} = 0\) 对 \(i\in \mathrm{I - J}\) 。如果 \(\mathfrak{a}_{\beta}\) 表示按第 2 号开头（IV，第26页）所定义的 A{[}{[}I{]}{]} 的理想，那么

\[
u \in \mathrm{A} \left[ \left(\mathrm{X} _ {i}\right) _ {i \in I} \right] \cap \mathfrak {a} _ {\beta} \Rightarrow \psi (u) \in \mathrm{V}.
\]

这表明 \(\psi\) 是连续的，如果我们给 \(\mathbf{A}[(\mathbf{X}_i)_{i\in I}]\) 赋予由 \(\mathbf{A}[[I]]\) 的拓扑诱导的拓扑。由于 \(\mathbf{E}\) 是分离的且完备的，\(\psi\) 延拓为一个连续保单位同态 \(\varphi\) ，它从 \(\mathbf{A}[[I]]\) 到 \(\mathbf{E}\) 中。我们有 \(\varphi (\mathbf{X}_i) = \psi (\mathbf{X}_i) = x_i\) 对所有 \(i\in I\) 。最后，设 \(\varphi^{\prime}\) 是从 \(\mathbf{A}[[I]]\) 到 \(\mathbf{E}\) 中的连续保单位同态，使得 \(\varphi^{\prime}(\mathbf{X}_{i}) = x_{i}\) 。我们有 \(\varphi^{\prime}(u) = \varphi (u)\) 对所有 \(u\in \mathbf{A}[(\mathbf{X}_i)_{i\in I}]\) ，因此 \(\varphi^{\prime} = \varphi\) ，因为 \(\mathbf{A}[(\mathbf{X}_i)_{i\in I}]\) 在 \(\mathbf{A}[[I]]\) 中稠密。

我们沿用先前的记号。若 \(u \in \mathbf{A}[[\mathbf{I}]]\) ，\(u\) 在 \(\varphi\) 下的像记作 \(u(\mathbf{x})\) 或 \(u((x_i)_{i \in \mathbf{I}})\) （或亦记作 \(u(x_1, \ldots, x_n)\) ，若 \(\mathbf{I} = \{1, 2, \ldots, n\}\) ），它称为 \(\mathbf{E}\) 中以 \(x_i\) 代替 \(X_i\) 代入 \(u\) 所得的元素，或 \(u\) 当 \(x_i\) 是 \(X_i\) 的值时的值；也称为 \(u\) 当 \(X_i = x_i\) 时的值。特别地，我们有 \(u = u((X_i)_{i \in \mathbf{I}})\) 。

设 \(E'\) 是一个结合、交换、含单位元、线性拓扑化、分离且完备的 A-代数。设 \(\lambda\) 是从 E 到 \(E'\) 中的连续保单位同态，且 \((x_i)_{i \in I}\) 是 E 中满足命题 4（IV，第28页）的条件 \(a)\) 与 \(b)\) 的元素族。族 \((\lambda(x_i))_{i \in I}\) 满足同样的条件 \(a)\) 与 \(b)\) 。对每个 \(u \in A[[I]]\) 我们有

\[
\lambda \left(u \left(\left(x _ {i}\right) _ {i \in I}\right)\right) = u \left(\left(\lambda \left(x _ {i}\right)\right) _ {i \in I}\right),\tag{3}
\]

因为映射 \(u \mapsto \lambda(u((x_i)_{i \in I}))\) 是从 \(A[[I]]\) 到 \(E'\) 中的连续保单位同态，它把 \(X_i\) 变为 \(\lambda(x_i)\) 对所有 \(i \in I\) 。

如果 \(J\) 与 \(K\) 是两个集合，我们用 \(A_{J,K}\) 表示满足下列条件的所有族 \((g_j)_{j \in J}\) 构成的集合：

\begin{enumerate}
\def\labelenumi{(\roman{enumi})}
\item
  对所有 \(j \in \mathbf{J}\) ，\(g_j\) 是 \(\mathbf{A}[[\mathbf{K}]]\) 中常数项为幂零元的元素；
\item
  若 \(J\) 是无限集，则 \(g_{j}\) 沿 \(J\) 的有限子集的补集所构成的滤子趋于 0 。
\end{enumerate}

我们注意到，若 J 有限，则每个族 \((g_{j})_{j\in J}\) （由 A{[}{[}K{]}{]} 中无常数项的形式幂级数组成）都属于 \(A_{J,K}\) 。

设 \((g_j)_{j\in J}\) 属于 \(\mathbf{A}_{\mathrm{J},\mathbf{K}}\) 。由命题 3 的推论（IV，第28页），我们有 \(\lim_{n\to \infty}g_j^n = 0\) 对所有 \(j\in J\) 。设 \(f\in \mathbf{A}[[\mathrm{J}]]\) ；我们可以把 \(g_{j}\) 代入

 \(j\) 在 \(f\) 中所对应的变量，并得到一个形式幂级数 \(f((g_j)_{j \in J})\) ，它属于 A{[}{[}K{]}{]} 。此外，映射 \(f \mapsto f((g_j)_{j \in J})\) 是从 A-代数 A{[}{[}J{]}{]} 到 A{[}{[}K{]}{]} 的连续同态。

特别地，若 \(J = \{1, \ldots, p\}\) 且 \(f \in A[[X_1, \ldots, X_p]]\) ，我们可以对每个 \(X_j\) 代入一个无常数项的形式幂级数 \(g_j \in A[[K]]\) ；此代入的结果记作 \(f(g_1, \ldots, g_p)\) 。

设 \(\mathbf{x} = (x_k)_{k\in K}\) 是 E 中满足命题 4（IV，第28页）的条件 \(a\) 与 \(b\) ) 的元素族。让我们应用 (3)，取 \(\lambda\) 为由 \(u\mapsto u(\mathbf{x})\) 给出的、把 A{[}{[}K{]}{]} 映入 E 的同态；我们得到

\[
f \left(\left(g _ {j}\right) _ {j \in J}\right) (\mathbf {x}) = f \left(\left(g _ {j} (\mathbf {x})\right) _ {j \in J}\right).\tag{4}
\]

设 \(\mathbf{f} = (f_i)_{i\in I}\in (\mathbf{A}[[J]])^I\) 与 \(g = (g_j)_{j\in J}\in \mathbf{A}_{J,K}\) 。我们用 \(\mathbf{f}(\mathbf{g})\) 或 \(\mathbf{f}\circ \mathbf{g}\) 表示元素 \((f_{i}((g_{j})_{j\in J}))_{i\in I}\) ，它属于 \((\mathbf{A}[[K]])^I\) 。若 \(\mathbf{f}\in \mathbf{A}_{I,J}\) ，则 \(\mathbf{f}\circ \mathbf{g}\in \mathbf{A}_{I,K}\) ，因为映射 \(f\mapsto f((g_j)_{j\in J})\) （从 \(\mathbf{A}[[I]]\) 到 \(\mathbf{A}[[K]]\) ）是连续的。

设 \(\mathbf{f} \in (\mathbf{A}[[J]])^{\mathrm{I}}\) ，\(\mathbf{g} \in \mathbf{A}_{\mathrm{J},\mathrm{K}}\) ，\(\mathbf{h} \in \mathbf{A}_{\mathrm{K},\mathrm{L}}\) 。则 \(\mathbf{g} \circ \mathbf{h} \in \mathbf{A}_{\mathrm{J},\mathrm{L}}\) ，并且由 (4)，我们有

\[
(\mathbf {f} \circ \mathbf {g}) \circ \mathbf {h} = \mathbf {f} \circ (\mathbf {g} \circ \mathbf {h}).\tag{5}
\]

